
\textbf{Lemme 1.7.} \uline{Soient} @@M0@@ \uline{un anneau principal}, @@M1@@. \uline{Alors l'homomorphisme canonique} :
@@M2@@
\uline{est un isomorphisme}.

Comme @@M3@@ est régulier, l'homomorphisme @@M4@@ est un isomorphisme. Donc @@M5@@ est engendré par les classes des @@M6@@-Modules localement libres de type fini, i.e.\ par les classes des @@M7@@-modules projectifs de type fini. Comme @@M8@@ est principal, tout module projectif de type fini est libre, d'où la surjectivité. L'injectivité vient de ce que @@M9@@ est @@M10@@-augmenté.

\textbf{Proposition 1.8.} \uline{Soient} @@M11@@ \uline{un schéma noethérien}, @@M12@@ \uline{un} @@M13@@\uline{-Module localement libre de rang} @@M14@@, @@M15@@ \uline{le fibré vectoriel associé à} @@M16@@, @@M17@@ \uline{le morphisme structural de} @@M18@@, \uline{et} @@M19@@ \uline{la section nulle de} @@M20@@. \uline{Alors l'homomorphisme composé} :
@@M21@@
(\uline{où} @@M22@@ \uline{est défini d'après} IV 2.12.1 \uline{et} 1.5) \uline{est la multiplication par l'élément} @@M23@@ \uline{de} @@M24@@, \uline{en notant} @@M25@@ \uline{la classe de} @@M26@@ \uline{dans} @@M27@@, \uline{et en posant}
@@M28@@

Cet énoncé est l'analogue, pour le @@M29@@, de VII 2.7.

Soit @@M30@@ un @@M31@@-Module cohérent, de classe @@M32@@ dans @@M33@@. Alors @@M34@@ est la classe dans @@M35@@ du complexe @@M36@@, donc de @@M37@@ considéré comme @@M38@@-Module par l'intermédiaire de @@M39@@. Pour calculer @@M40@@, il faut donc prendre la classe de @@M41@@.

